\section{Resource-Constrained Control Design}
\label{app:resource_constrained_control}

The theory in Section~\ref{sec:feedback_efficiency} separates three questions that are often conflated in controller design: how strongly the population can be moved, how coherently controller influence is converted into that motion, and how efficiently the finite-time sensing--actuation budget is converted into directed current. This appendix turns susceptibility, the two bounded efficiencies, and an explicit operational interaction cost into a practical design prescription.

\subsection{Minimum-cost control at prescribed response}
\label{app:min_cost_control}

Let
\[
\boldsymbol{\varphi}=(q_c,b,\theta,\beta,\ldots)
\]
denote the tunable feedback-controller configuration, and let \(\rho(n)\) be the population-state occupancy over the evaluation window. The microscopic channel parameters \(h\) and \(\gamma\) are treated separately: unless the implementation exposes a mechanism that changes them directly, they are calibrated properties of the population response rather than free design knobs.

Since the state-local susceptibility is defined conditionally on advocacy, the expected controller-induced response per round is
\begin{equation}
R(\boldsymbol{\varphi};h,\gamma)
=
\sum_{n=0}^{N}
\rho(n)\,
a_n(\boldsymbol{\varphi})\,
\chi_{h,\gamma}\!\left(\frac{n}{N}\right).
\label{eq:app_expected_response}
\end{equation}
A direct operational interaction cost is
\begin{equation}
C_{\mathrm{op}}(\boldsymbol{\varphi})
=
q_c
+
b\sum_{n=0}^{N}\rho(n)a_n(\boldsymbol{\varphi}),
\label{eq:app_operational_cost}
\end{equation}
where the first term counts the agents sensed every round and the second counts the expected number of controlled microscopic opportunities.

The most direct controller-design problem is therefore not to maximize an efficiency in isolation, but to achieve a prescribed behavioral response at minimum operational cost,
\begin{equation}
\boxed{
\boldsymbol{\varphi}^\star
=
\arg\min_{\boldsymbol{\varphi}} C_{\mathrm{op}}(\boldsymbol{\varphi})
\qquad
\text{subject to}
\qquad
R(\boldsymbol{\varphi};h,\gamma)\ge R_{\mathrm{req}} .}
\label{eq:app_min_cost_problem}
\end{equation}
Equivalently, for a fixed interaction budget one may maximize the achievable response,
\begin{equation}
\boxed{
\boldsymbol{\varphi}^\star
=
\arg\max_{\boldsymbol{\varphi}} R(\boldsymbol{\varphi};h,\gamma)
\qquad
\text{subject to}
\qquad
C_{\mathrm{op}}(\boldsymbol{\varphi})\le C_{\max}.}
\label{eq:app_max_response_problem}
\end{equation}
These formulations make the practical question explicit: once the required response has been fixed, how much sensing and actuation can be removed before performance falls below the desired level?

For the exact finite-compliance \(q=1\) reference,
\begin{equation}
\chi_{h,\gamma}(x)
=
[\sigma(h)-x]\Lambda_{b,\gamma},
\qquad
\Lambda_{b,\gamma}
=
1-\left(1-\frac{\gamma}{N}\right)^b .
\label{eq:app_susceptibility_recall}
\end{equation}
The marginal gain from one additional controlled opportunity is
\begin{equation}
\Lambda_{b+1,\gamma}-\Lambda_{b,\gamma}
=
\frac{\gamma}{N}
\left(1-\frac{\gamma}{N}\right)^b,
\label{eq:app_marginal_actuation_gain}
\end{equation}
which decreases monotonically with \(b\). Thus stronger actuation continues to increase the mean response toward the set point, but each additional controlled opportunity contributes less than the previous one. The same logic applies to sensing: increasing \(q_c\) is useful only insofar as the improved state estimate changes the policy in a way that improves the required response.

\subsection{How the two efficiencies guide feasible designs}
\label{app:efficiency_guided_design}

The operational problem above specifies how to purchase less control. The two
bounded efficiencies from Section~\ref{sec:feedback_efficiency} then determine
whether a feasible controller uses its influence and thermodynamic budget well.

The information--response efficiency is
\begin{equation}
\eta_{\mathrm{IR}}(n)
=
\frac{2a_n(1-a_n)\,\chi_{h,\gamma}(n/N)^2}
{(\ln 2)\,T_\pi(n)}
\le 1,
\label{eq:app_eta_ir}
\end{equation}
with $T_\pi(n)=I(U_k;n_{k+1}\mid n_k=n)$. It measures whether the
action-dependent change in the full next-state distribution is expressed as a
coherent shift in mean target fraction. Among settings with comparable cost and
response, one should prefer regimes in which $\eta_{\mathrm{IR}}$ remains high
over the states actually visited.

The thermodynamic efficiency is
\begin{equation}
\boxed{
\eta_{\mathrm{th},k}
=
\frac{hJ_{c,k}}{\Sigma_k-\Delta S_{\mathrm{sys},k}}
=
\frac{hJ_{c,k}}{hJ_{c,k}+I(n_k;Y_k)}
\le 1.}
\label{eq:app_eta_th}
\end{equation}
Here $\Delta S_{\mathrm{sys},k}$ is the transient storage term and
$\Sigma_k-\Delta S_{\mathrm{sys},k}$ is the non-storage finite-time control
expenditure. The efficiency therefore reports the fraction of the generalized
sensing--actuation expenditure expressed as affinity-weighted directed current.
Its derivation from the forward/reverse path ratio is given in
Appendix~\ref{app:thermodynamic_efficiency_derivation}.

Two practical rules follow directly. First, more sensing is not automatically
better: if increasing $q_c$ raises $I(n_k;Y_k)$ without a proportional increase
in $hJ_c$, it lowers $\eta_{\mathrm{th}}$. Second, more actuation is useful only
while it continues to raise the required response; the exact susceptibility
already shows diminishing marginal gains through
$\Lambda_{b+1,\gamma}-\Lambda_{b,\gamma}$.

Taken together, the design prescription is
\begin{equation}
\boxed{
\begin{aligned}
&\text{(i) specify the minimum acceptable response or current},\\
&\text{(ii) minimize operational cost subject to that requirement},\\
&\text{(iii) use }\eta_{\mathrm{IR}}\text{ to reject incoherent action channels},\\
&\text{(iv) among feasible controllers, prefer larger }\eta_{\mathrm{th}},\\
&\text{(v) retain }h\text{ and }\gamma\text{ as separate directional and kinetic diagnostics}.
\end{aligned}}
\label{eq:app_design_recipe}
\end{equation}
Thus susceptibility determines where control has leverage, the operational cost
determines how much control must be purchased, $\eta_{\mathrm{IR}}$ determines
whether controller influence becomes coherent mean motion, and
$\eta_{\mathrm{th}}$ determines whether sensing and actuation are converted
efficiently into directed current.

\subsection{Finite-time and nonstationary interpretation}
\label{app:nonstationary_design}

No stationarity assumption is required in the construction above. Let
\begin{equation}
\rho_k(n)=\Pr(n_k=n)
\end{equation}
denote the transient population distribution at round \(k\). Over a finite horizon of \(H\) rounds, the occupancy entering Eqs.~(\ref{eq:app_expected_response})--(\ref{eq:app_operational_cost}) can be interpreted as the finite-time visitation distribution
\begin{equation}
\bar{\rho}_H(n)
=
\frac{1}{H}
\sum_{k=0}^{H-1}\rho_k(n).
\label{eq:app_finite_time_occupancy}
\end{equation}
Hence the operational design problem averages over the transient states actually visited during the experiment rather than over an equilibrium distribution. Likewise, Eq.~(\ref{eq:app_th_second_law}) retains the finite-time boundary contribution \(\Delta S_{\mathrm{sys},k}\), which accounts for the fact that the population ensemble changes during the feedback cycle. A stationary distribution is therefore only a special long-time limit of the present formulation, not an assumption required by it.
